\section{Practical method, quality control and limitations}

\subsection{Choosing a source image}

Percepta can process any conventional image, but not every source produces an equally convincing illusion. The encoder has finite spatial bandwidth: details smaller than the pattern period cannot be represented reliably. Sources with a clear hierarchy therefore perform best.

Favour images with:

\begin{itemize}
  \item one dominant subject occupying a substantial fraction of the frame;
  \item recognisable silhouettes and medium-scale features;
  \item adequate separation between subject and background;
  \item controlled highlights and shadows;
  \item colour variation that remains meaningful after local averaging.
\end{itemize}

Potentially difficult sources include distant crowds, foliage, text-heavy scenes, repeating fabrics, fine hair against a busy background and images whose identity depends on one-pixel-scale detail.

\subsection{Pattern selection by subject}

\begin{table}[H]
\centering
\caption{Starting pattern recommendations.}
\begin{tabularx}{\textwidth}{>{\bfseries}l >{\raggedright\arraybackslash}X >{\raggedright\arraybackslash}X}
\toprule
Subject & Recommended starting modes & Reason \\
\midrule
Portrait & Vertical stripes, Halftone, Spiral & Faces tolerate controlled simplification and are recognised from medium-scale features \\
Landscape & Horizontal stripes, Hexagonal halftone & Horizons and broad tonal regions align naturally with lateral or isotropic sampling \\
Architecture & Diagonal stripes, Hexagonal halftone & Diagonal carriers avoid simple axis coincidence; hex grids reduce directional bias \\
Graphic logo & Vertical/Horizontal stripes, Halftone & Flat regions and strong boundaries translate clearly into coverage \\
Highly textured image & Dense halftone or dense stripes & More samples are required to preserve local variation \\
Abstract source & Spiral or strong Colour separation & Pattern character can take precedence over literal reconstruction \\
\bottomrule
\end{tabularx}
\end{table}

These are starting points rather than rules. A deliberate mismatch between subject and carrier can be aesthetically productive.

\subsection{A controlled optimisation sequence}

Changing several controls at once makes it difficult to identify the cause of improvement. The following one-variable sequence is recommended.

\begin{enumerate}
  \item \textbf{Lock framing.} Place the subject and do not change the crop during the first comparison.
  \item \textbf{Compare pattern families at their defaults.} Choose the structure whose close-range appearance suits the project.
  \item \textbf{Sweep density.} Test at least one sparse, one reference and one dense value.
  \item \textbf{Set Strength.} Increase until important tonal regions separate, then step back if highlights merge.
  \item \textbf{Set Source contrast.} Use the smallest increase that restores recognisability.
  \item \textbf{Add Colour separation.} Stop before the channels become independent copies at the intended viewing distance.
  \item \textbf{Tune pattern-specific smoothing or minimum size.}
  \item \textbf{Test final dimensions.} Increase output size and check adaptive values.
  \item \textbf{Proof the export.} Inspect actual PNG/TIFF/SVG/PDF output, not only the live preview.
\end{enumerate}

\subsection{Quantitative quality indicators}

Perceptual judgement remains central, but several numerical indicators can help diagnose a result.

\subsubsection{Low-pass reconstruction error}

Let $\mathcal{L}$ be a low-pass operator approximating distant viewing. Compare the filtered output with a filtered version of the framed source:

\begin{equation}
E_{\mathrm{LP}}=
\frac{1}{3WH}
\sum_{x,y}
\left\|
\mathcal{L}[\mathbf{I}_{\mathrm{out}}](x,y)
-
\mathcal{L}[\mathbf{C}'](x,y)
\right\|_2^2.
\label{eq:lperror}
\end{equation}

A smaller value indicates better preservation at the selected integration scale. It does not measure the artistic quality of the visible pattern.

\subsubsection{Structural carrier visibility}

A high-pass operator $\mathcal{H}$ can estimate how strongly the geometric pattern remains visible:

\begin{equation}
V_{\mathrm{carrier}}=
\frac{1}{3WH}
\sum_{x,y}
\left\|\mathcal{H}[\mathbf{I}_{\mathrm{out}}](x,y)\right\|_2^2.
\label{eq:carriervis}
\end{equation}

Percepta's artistic objective is not to minimise $V_{\mathrm{carrier}}$. A successful result often balances low reconstruction error with enough high-frequency energy to make the construction visually compelling.

\subsubsection{Gamut and clipping rate}

The clipping fraction in Equation~\eqref{eq:clipfraction} can be extended to exported RGB values. Large areas at exact 0 or 1 may be intended, but unexpected clipping can reveal excessive contrast or overlapping channel geometry.

\subsection{Troubleshooting guide}

\begin{tabularx}{\textwidth}{>{\bfseries}p{0.23\textwidth} >{\raggedright\arraybackslash}X >{\raggedright\arraybackslash}X}
\toprule
Symptom & Likely cause & Corrective action \\
\midrule
Source is not recognisable & Pattern too sparse; source too small in frame; weak tonal separation & Increase stripe count or reduce spacing; crop tighter; raise contrast moderately \\
Image becomes nearly white & Excessive Strength, contrast or overlap & Lower Strength; reduce contrast; reduce minimum size; inspect bright-source clipping \\
RGB copies are visible separately & Colour separation too large for density/viewing distance & Reduce separation; use denser geometry; increase viewing distance \\
Dark regions contain too much colour & Minimum geometry too large or smoothing spreads bright features & Reduce minimum size; reduce smoothing; inspect source black level \\
Thin lines flicker or break & Insufficient raster resolution or antialiasing & Increase output size; enable smoothing; export vector when appropriate \\
Diagonal edges look jagged & Raster grid aliasing & Increase supersampling/output size; use antialiased vector paths \\
Spiral turns collide & Turn spacing too small relative to width/separation & Increase turn spacing; lower Strength or separation \\
Animation pauses at loop point & Duplicate terminal frame or endpoint easing & Omit duplicated endpoint; adjust loop timing; use ping-pong phase consistently \\
Export is unexpectedly large & Dense vector geometry or many animation frames & Use lossless raster; reduce dimensions/FPS/duration; stream frames to encoder \\
Print loses coloured gaps & Printer resolution, dot gain or scaling & Make a final-size proof; increase feature width; avoid driver resampling \\
\bottomrule
\end{tabularx}

\subsection{Computational complexity}

For a raster output of $W\times H$ pixels, preprocessing and mask composition are generally $\mathcal{O}(WH)$. Stripe renderers can also be implemented in $\mathcal{O}(WH)$ by evaluating carrier phase per pixel or by rasterising paths whose total sampled length scales with image area at fixed density.

A halftone with spacing $P$ contains approximately

\begin{equation}
N_{\mathrm{dots}}\approx\frac{WH}{P^2}
\label{eq:dotcount}
\end{equation}

for a square grid, and

\begin{equation}
N_{\mathrm{hex}}\approx
\frac{2WH}{\sqrt{3}P^2}
\label{eq:hexcount}
\end{equation}

for a triangular lattice. Vector export cost grows with the number of primitives, whereas raster cost remains dominated by the output pixel count.

For an Archimedean spiral extending to radius $R_{\max}$ with turn spacing $T$, the number of turns is approximately

\begin{equation}
N_{\mathrm{turns}}\approx\frac{R_{\max}}{T}.
\end{equation}

The arc length from $\theta=0$ to $\theta=\Theta$ is

\begin{equation}
\ell(\Theta)=
\int_0^{\Theta}
\sqrt{r(\theta)^2+\left(\frac{\mathrm{d}r}{\mathrm{d}\theta}\right)^2}
\,\mathrm{d}\theta.
\label{eq:spirallength}
\end{equation}

Dense spirals therefore require many path samples and can be more expensive than simple stripes.

\subsection{Known limitations}

\begin{warningbox}
Percepta is an artistic perceptual encoder, not a reversible compression method. Fine detail is intentionally discarded, and the original image cannot be recovered exactly from the output.
\end{warningbox}

The principal limitations are:

\begin{itemize}
  \item \textbf{Source dependence.} Optimal settings vary with composition, contrast and scale.
  \item \textbf{Observer dependence.} Visual acuity, colour vision and viewing conditions alter the effect.
  \item \textbf{Device dependence.} Displays and printers differ in gamut, resolution, transfer function and scaling.
  \item \textbf{Sampling loss.} Features smaller than the carrier period cannot be represented faithfully.
  \item \textbf{Edge interaction.} RGB separation can create coloured ghosts around high-contrast boundaries.
  \item \textbf{Raster/vector differences.} Antialiasing and path tessellation can cause small differences between preview and export.
  \item \textbf{Animation cost.} High-resolution sequences multiply rendering and encoding requirements.
\end{itemize}

\subsection{Implementation and distribution}

The current application is designed around a compact Python/PyQt6 project. A clean distribution can contain the executable source entry point and an assets folder, with the launcher responsible for component checks and packaged execution.

The runtime requirements used by the project include Python 3.10 or newer, PyQt6, NumPy, Pillow, imageio and imageio-ffmpeg. A packaged Windows build can use PyInstaller. The startup sequence may display the Percepta splash screen, check components once per application version and store state in the user's application-data directory.

Still-image export should be deterministic for a fixed source, framing state, pattern and parameter set. Reproducibility is improved by storing:

\begin{itemize}
  \item application version;
  \item source file identity and dimensions;
  \item crop and affine framing parameters;
  \item pattern name and every common/pattern-specific parameter;
  \item output dimensions and export format;
  \item animation timing and easing when applicable.
\end{itemize}

A future preset or sidecar format could serialise these values in JSON.

\subsection{Possible extensions}

The existing model naturally supports several extensions without changing its core identity:

\begin{itemize}
  \item perceptual or device-calibrated coverage functions $F_q^{-1}$;
  \item automatic density selection based on source spatial frequency;
  \item GPU rasterisation for high-resolution animation;
  \item batch rendering and preset comparison sheets;
  \item objective low-pass reconstruction previews and error maps;
  \item additional continuous paths or lattice families;
  \item colour-management metadata and printer-specific proof modes.
\end{itemize}

These are development possibilities rather than current user features.


